\chapter{Current Validation and Observations}
\label{chap:validation}

This chapter reports two different kinds of evidence, and keeps them
clearly apart. The first is an automated test suite that checks each
component's logic in isolation, against mocked external systems. The
second is a small, real-system pilot that checks whether the same
components behave correctly against the real PyPI registry, a real Ollama
model, a real Docker daemon, real GitHub repositories, and, for the
result-logging and knowledge-graph components, the real Morph-KGC mapping
engine and the FAIR Jupyter knowledge graph's own source data.

The pilots reported through most of this chapter each
exercise one component's own stages against its relevant real external
system in isolation. \Cref{sec:v-i7} reports a further pilot that
drives several components together, unattended, in one run via
the orchestrator (\Cref{sec:orchestrator}), still over a
small number of records. \Cref{sec:v-final-eval} then reports the thesis's
final evaluation: the same orchestrator, unchanged, run once over the
complete 187-row evaluation split described in \Cref{chap:problem-dataset}.

\section{Automated Testing}
\label{sec:v-testing}

At the implementation state this thesis reports, the repository's test
suite contains \textbf{332 passing tests}, organised by the component they
exercise. This is a measure of software correctness, not of repair quality
or explanation quality: it confirms that each component's own logic behaves
as designed against controlled inputs, not how well the system performs on
real notebooks at scale.

\begin{table}[h]
  \centering
  \begin{tabular}{|p{5.2cm}|r|p{5.8cm}|}
    \hline
    \textbf{Area} & \textbf{Tests} & \textbf{What is checked} \\
    \hline
    Context extraction and classification & 9 &
      Root-cause hints, confidence levels, and the legacy-hint
      acceptance/rejection logic from \Cref{sec:m-classification}. \\
    \hline
    Explanation schema, prompt rendering, and runner & 29 &
      Prompt template rendering, JSON schema validation, and the
      retry-on-failure logic from \Cref{sec:m-explanation}. \\
    \hline
    PyPI retrieval and compatibility evidence & 120 &
      Import-name mapping, candidate-version filtering, Python-compatibility
      checks, and the compatibility-evidence registry lookup from
      \Cref{sec:m-repair-retrieval,sec:m-repair-compat}. \\
    \hline
    Repair proposal validation and agent orchestration & 78 &
      Grounding validation, argument-list construction, and rejection of
      an invented name or an out-of-range version, from \Cref{sec:m-repair}. \\
    \hline
    Fix application, Docker orchestration, and outcome classification & 74 &
      Command re-validation, the dataset re-join, outcome classification,
      and cleanup behaviour, from \Cref{sec:m-fixapply}. \\
    \hline
    ResultLogger joining and SQLite persistence & 15 &
      The classification/explanation join, the positional
      repair-proposal/fix-application join and its duplicate-attempt
      regression case, the missing-outcome run\_id fallback, and
      end-to-end logging runs against temporary fixtures, from
      \Cref{sec:m-resultlogger-join,sec:m-resultlogger-row}. \\
    \hline
    Repair-attempt CSV export & 7 &
      \texttt{notebook\_id} resolution via the classification-dataset join,
      CSV column order, \texttt{NULL}-to-empty-field conversion, and JSON
      blob round-tripping, from \Cref{sec:m-resultlogger-csv}. \\
    \hline
    \textbf{Total} & \textbf{332} & \\
    \hline
  \end{tabular}
  \caption{Automated tests by component area, at the current implementation state.}
  \label{tab:v-tests}
\end{table}

Dataset preparation has no dedicated automated test file of its own;
instead, its output is checked by independently recomputing the counts in
\Cref{chap:problem-dataset} directly from its output files and confirming
they agree exactly, which they do.

Across these areas, the tests specifically exercise: rejection of malformed
or schema-invalid model output; rejection of a repair proposal that invents
a package name or a version outside the retrieved evidence; the shell-safety
checks described in \Cref{sec:m-fixapply}, including a static check that
neither \texttt{shell=True}, \texttt{os.system}, \texttt{eval}, nor
\texttt{exec} appears anywhere in the components that build or run a
command; the three-outcome classification logic, including the conservative
same-error check; cleanup behaviour under both success and failure; and, for
ResultLogger, rejection of an ambiguous duplicate explanation record and
correct separation of two repair attempts sharing the same
\texttt{notebook\_execution\_id} into two distinct rows. All
of this runs against mocked subprocesses, a mocked Ollama endpoint, and
mocked HTTP responses. This is deliberate and also a limitation: a mocked
test can prove a component's logic is internally consistent, but, as
\Cref{sec:m-fixapply-bugs} already showed for three real bugs it could not
have caught, it cannot prove the component behaves correctly against a real
operating system, a real Docker daemon, or a real network. That is exactly
what the pilots below were run to check.

\section{RAGRepairAgent Implementation Validation}
\label{sec:v-i4}

\subsection{Retrieval Against Real PyPI}

The five import names used in the original retrieval-design check
(\texttt{sklearn}, \texttt{umap}, \texttt{pkg\_resources}, \texttt{scipy},
\texttt{dms\_variants}) were re-checked against the real, production
\texttt{retrieve()} entry point and the live PyPI registry. All four mapped
names resolved correctly to their real PyPI distributions with plausible,
current candidate versions; \texttt{dms\_variants}, which has no verified
mapping, correctly returned \texttt{mapping\_unknown} with no PyPI request
made at all, confirming under real conditions -- not only in a mocked test
-- that an unmapped import never reaches the network.

A later, separate check confirmed a fix to the mapping table itself.
\texttt{cv2} and \texttt{skimage} are both common, usable failing modules in
the dataset (\Cref{tab:pd-top-modules}) that had no entry in
\texttt{config/package\_mapping.yaml}, so the retriever returned
\texttt{mapping\_unknown} for both instead of resolving them. Adding
\texttt{cv2: opencv-python} and \texttt{skimage: scikit-image} was checked
against the real PyPI registry, alongside a regression check confirming
that the existing \texttt{sklearn} mapping still resolved correctly and was
not disturbed by the addition.

\subsection{Repair Proposals Against a Real Model}

A small pilot ran \texttt{rag\_repair\_agent.py}'s real entry point, using a
real local Ollama model, over four records:

\begin{itemize}
  \item An excluded, \texttt{system\_library} record correctly abstained
        with zero PyPI or Ollama calls, confirming the eligibility gate
        works under real conditions.
  \item A \texttt{missing\_package} record (\texttt{sklearn}) timed out
        twice under the originally configured 120-second ceiling, then
        succeeded on the first attempt once the ceiling was raised to 300
        seconds -- the direct evidence behind the timeout change reported
        in \Cref{sec:m-repair}.
  \item The same record, rerun at 300 seconds, produced a schema-valid,
        grounded proposal: install \texttt{scikit-learn}.
  \item A \texttt{wrong\_version} record (the \texttt{scipy.integrate.cumtrapz}
        case from \Cref{tab:m-compat}, traced in full in
        \Cref{sec:m-worked-example}) produced a grounded
        \texttt{pin\_version} proposal for \texttt{scipy==1.13.1}, using the
        real compatibility-evidence intersection on the first attempt.
\end{itemize}

\subsection{What This Pilot Does and Does Not Show}

It shows that real PyPI retrieval, real import-name mapping, and real model
proposal generation operate correctly across RAGRepairAgent's own
retrieval-to-proposal flow, and that timeout and retry behaviour degrades
safely rather than hanging or fabricating a result. It does not show an
installed package, a modified environment, or a re-executed notebook --
none of that happened at this step, by design; validating a proposal is
FixApplicator's job, covered next.

\section{FixApplicator Implementation Validation}
\label{sec:v-i5}

A second pilot ran \texttt{apply\_and\_validate()} against a real Docker
daemon and real GitHub repositories, reusing already-validated repair
proposals from the pilot above. The upstream sibling database was found
locked by another process during this pilot, which meant commit metadata
for one lookup had to fall back to its documented "unavailable" behaviour --
itself a real, unplanned exercise of that fallback path, not a simulated
one.

\begin{table}[h]
  \centering
  \small
  \begin{tabular}{|p{1.7cm}|p{3.2cm}|p{2.6cm}|p{6.0cm}|}
    \hline
    \textbf{Case} & \textbf{Original error} & \textbf{Fix applied} & \textbf{Result} \\
    \hline
    1 -- \texttt{sklearn} & \texttt{ModuleNot\-FoundError}: sklearn & install \texttt{scikit-learn} &
      Original error gone; \texttt{still\_failing} on a later, unrelated \texttt{No module named 'pandas'}. \\
    \hline
    2 -- \texttt{scipy} & \texttt{ImportError}: cumtrapz & pin \texttt{scipy==1.13.1} &
      Original error gone; \texttt{still\_failing} on a later, unrelated \texttt{No module named 'tf\_keras'}. \\
    \hline
    3 -- controlled failure & synthetic: non-existent repository & n/a &
      \texttt{apply\_error}, failure stage \texttt{clone}; no Docker artefacts created. \\
    \hline
    4 -- \texttt{cv2} & \texttt{ModuleNot\-FoundError}: cv2 & install \texttt{opencv-python} &
      Original error gone; \texttt{still\_failing} on a missing \emph{system} library (\texttt{libxcb.so.1}), a known characteristic of \texttt{opencv-python} on a minimal base image. \\
    \hline
    5 -- \texttt{sklearn} (2nd repo) & \texttt{ModuleNot\-FoundError}: sklearn & install \texttt{scikit-learn} &
      Original error gone; \texttt{still\_failing} on a later, unrelated \texttt{No module named 'tensorflow'}. \\
    \hline
  \end{tabular}
  \caption{Real-Docker pilot cases. Case 3 is a deliberately synthetic
  infrastructure-failure test, not a dataset row.}
  \label{tab:v-i5-cases}
\end{table}

Cases 4 and 5 were chosen deliberately, after cases 1--3, specifically to
give the pilot the best possible chance of reaching a completely clean
\texttt{fixed} result: both repositories have no declared requirements
file, no \texttt{setup.py}, and exactly one notebook, minimising the chance
of a second, unrelated dependency gap.

\subsection{Key Observation: Fixing One Dependency Can Reveal Another}
\label{sec:v-key-observation}

Every one of the five real cases above -- spanning four different
repositories, two different fix actions, and packages with no shared code
lineage -- resolved the specific error it targeted and then failed again on
a different, previously invisible problem. This is not a coincidence and it
is not a defect in FixApplicator's own classification logic, which is
separately unit-tested and proven capable of reporting a clean
\texttt{fixed} result the moment one occurs. It is a property of the
dataset: the Docker pipeline's error classification records only the
\emph{first} import failure a notebook hits. A real research notebook can
easily depend on several packages that are each individually missing; once
the first one is repaired, execution proceeds further and can reach the
next one, which was simply invisible until that point. The conservative
\texttt{same\_as\_original\_error} check (\Cref{chap:methodology})
correctly distinguished "a different, new problem" from "the fix did not
work" in every one of these cases.

No completely clean \texttt{fixed} outcome
was produced by this pilot, despite it being deliberately designed to find
one. This is stated plainly rather than minimised: a clean outcome remains
achievable in principle on a single-dependency notebook. The two-round
repair mode that feeds a \texttt{still\_failing} result's new error
back into a second repair attempt was demonstrated for
real on exactly cases 1 and 2 above (\Cref{sec:v-i7}); in both real
attempts the second round correctly triggered and correctly abstained
rather than proposing anything ungrounded. \Cref{sec:v-final-eval} reports
that this same pattern -- a real repair removing its targeted error, but the
notebook remaining unfixed because a second, previously hidden problem
takes its place -- recurs throughout the full 187-row evaluation split, and
not only in this small pilot.

\section{Result Logging and Knowledge-Graph Implementation Validation}
\label{sec:v-i6}

\subsection{Notebook-Identity Alignment with the FAIR Jupyter Knowledge Graph}
\label{sec:v-i6-alignment}

Linking a repair attempt to the notebook it repairs, as the RML mapping in
\Cref{sec:m-resultlogger-rml} does, only works if the Docker pipeline's own
\texttt{notebook\_id} and \texttt{repository\_id} actually identify the same
notebooks and repositories as the existing FAIR Jupyter knowledge graph -- a
graph generated from the earlier, conda-based pipeline
(\Cref{sec:pd-why-docker}), not the Docker pipeline this thesis builds on.
This was checked directly rather than assumed. All \textbf{214} rows of the
classification dataset (\Cref{sec:m-classification}) were cross-checked
against the FAIR Jupyter knowledge graph's own \texttt{repositories.csv} and
\texttt{notebooks.csv} source files, comparing each row's resolved
repository path and notebook filename for an exact string match. Every one
of the 214 rows matched on both \texttt{repository\_id}$\to$\texttt{repository}
and \texttt{notebook\_id}$\to$\texttt{name}, with zero missing and zero
mismatched records.

Within the dataset inspected for this check, the Docker pipeline and the
FAIR Jupyter knowledge graph therefore use the same
\texttt{repository}/\texttt{notebook} primary keys for the same notebooks,
and a repair attempt can link directly to its existing notebook resource via
\texttt{notebook\_id} (\Cref{sec:m-resultlogger-identity}); no separate
crosswalk table between the two pipelines is required. This result
establishes alignment for the inspected FAIR Jupyter source dataset
specifically; it is not a general claim about every possible knowledge-graph
deployment built from a different data snapshot.

\subsection{Executing the RML Mapping with Morph-KGC}
\label{sec:v-i6-morphkgc}

The RML mapping described in \Cref{sec:m-resultlogger-rml} was not only
written; it was executed with Morph-KGC, the same mapping engine the
existing \texttt{run\_fairjupyter\_kg.sh} build script uses for every other
table in the knowledge graph. A \texttt{repair\_attempts.csv} file was
assembled from the real classification, explanation, and repair-proposal
pilot records already reported in \Cref{sec:v-i4}; since no live
fix-application batch output exists yet at full scale, a labelled synthetic
FixApplicator record stood in for one row, so that the outcome and
post-repair error fields could also be exercised in the generated RDF.

Running the mapping over this file produced valid Turtle, with no
RML/Turtle syntax error. The resulting RDF was then inspected directly
rather than merely trusted to have been generated correctly: a SQLite column
left \texttt{NULL} -- an abstained attempt's \texttt{install\_name}, for
instance -- produced no triple at all for the corresponding predicate, never
an empty-string literal; \texttt{repr:exception} and \texttt{repr:msg}
appeared only on the one row that actually carried a post-repair error; and
\texttt{prov:endedAtTime} was present, and correctly typed
\texttt{xsd:dateTime}, on every row.

The generated \texttt{repr:hadRepairAttempt} links were also checked against
real notebook resources, not merely against the CSV's own
\texttt{notebook\_id} column: the existing, unmodified
\texttt{notebooks.rml.ttl} was itself run against the real
\texttt{notebooks.csv}, producing 627{,}127 triples and confirming that this
file executes exactly as the rest of the knowledge graph expects, and the
notebook IRIs the repair mapping links to were confirmed present there,
typed \texttt{repr:Notebook} with the expected \texttt{pav:retrievedFrom}
link to their repository -- not merely matching by string inspection of two
separate CSV files. All validation output was written to a scratch
location; the FAIR Jupyter knowledge-graph checkout itself was confirmed
untouched throughout.

This validates that the mapping produces syntactically valid RDF for real
data, links correctly to existing notebook resources for the records
tested, and handles a missing value and the PROV-O time literal correctly.

It is, again, a pilot check against a small, real sample
assembled for this purpose, not a run over the complete dataset. The
\texttt{repair\_attempts} table itself has since been populated at the scale
of the full 187-row evaluation split (\Cref{sec:v-final-eval}); running this
same, unchanged RML mapping over that larger table to enrich the knowledge
graph accordingly has not been done as part of this evaluation, and
\Cref{chap:future-work} states what it still requires.

\section{Orchestrator Implementation Validation}
\label{sec:v-i7}

A further pilot ran the orchestrator itself
(\texttt{scripts/run\_pipeline.py}, \Cref{sec:orchestrator}) end to end --
classification through result logging, including the bounded two-round
repair loop -- against real Ollama, real PyPI, and a real Docker daemon,
never mocked, for two development-split records. Both reproduce, and
extend with a real Round 2 attempt, the cases already reported in
\Cref{sec:v-i5}.

\begin{table}[h]
  \centering
  \small
  \begin{tabular}{|p{1.3cm}|p{3.3cm}|p{3.3cm}|p{6.0cm}|}
    \hline
    \textbf{Case} & \textbf{Round 1} & \textbf{Round 2 trigger} & \textbf{Round 2 result} \\
    \hline
    A -- \texttt{sklearn} (id 8) &
      install \texttt{scikit-learn}; \texttt{still\_failing} on \texttt{No module named 'pandas'} &
      Fired: \texttt{pandas} reclassifies as \texttt{missing\_package}/usable &
      RAGRepairAgent called for real against \texttt{pandas}, abstained after two invalid-JSON
      responses -- the model hallucinated JSON echoing PyPI's own
      "unparseable filename" warnings for \texttt{pandas}'s many legacy Windows release
      artifacts, embedded in the prompt's evidence block. A genuine, unforced model
      failure mode; FixApplicator never reached Docker. \\
    \hline
    B -- \texttt{scipy}/cumtrapz (id 174) &
      pin \texttt{scipy==1.13.1}; \texttt{still\_failing} on \texttt{No module named 'tf\_keras'} &
      Fired: \texttt{tf\_keras} reclassifies as \texttt{missing\_package}/usable &
      RAGRepairAgent called for real against \texttt{tf\_keras}, abstained: PyPI retrieval
      returned \texttt{mapping\_unknown} because \texttt{tf\_keras} has no entry in
      \texttt{config/package\_mapping.yaml}. Confirms, rather than assumes, that this
      documented case cannot currently be repaired end to end; no mapping entry was added to
      force a different outcome. FixApplicator never reached Docker. \\
    \hline
  \end{tabular}
  \caption{Real orchestrator pilot: both cases exercise a genuine Round 2 trigger and a genuine,
  differently-caused RAGRepairAgent abstention, never a fabricated success.}
  \label{tab:v-i7-cases}
\end{table}

Round 1 of both cases reproduced the FixApplicator-only pilot results in
\Cref{tab:v-i5-cases} exactly, now driven by the orchestrator rather than
by hand, and with both rows carrying one shared \texttt{run\_id} across
both rounds -- confirmed directly in the resulting SQLite rows, not merely
assumed from the code. Each round produced its own, separate
\texttt{repair\_attempts} row (\Cref{sec:m-resultlogger}), never collapsed
into one.

Neither real Round 2 attempt in this small pilot produced a successful second
repair -- both triggered correctly and abstained correctly, for two
different, genuine reasons, rather than either being forced to succeed.
The controlled, mocked test
suite (\Cref{chap:methodology}) separately proves the second-round
mechanism itself -- trigger conditions, the cumulative environment, the
shared run identifier, and the two-round cap -- works correctly end to end
when RAGRepairAgent does propose and FixApplicator does apply a Round 2
fix. \Cref{sec:v-final-eval} reports how the second round behaved across
the full evaluation split: Round 2 fired for 48 notebooks, produced a real
second repair proposal for 12 of them, and, as in this pilot, none of those
12 second attempts produced a completely clean \texttt{fixed} result
either.

\section{Final Evaluation on the Reserved Evaluation Split}
\label{sec:v-final-eval}

Every result reported so far in this chapter comes from a small, hand-picked
pilot: a handful of records chosen to exercise one component, or the whole
orchestrator, against a real external system. This section reports the
opposite kind of evidence: the complete repair layer, unchanged from the
implementation described in \Cref{chap:methodology}, run once, unattended,
over every one of the 187 rows in the reserved \texttt{evaluation} split
(\Cref{chap:problem-dataset}), against a real Ollama model, the live PyPI
registry, and a real Docker daemon throughout. This is the thesis's final
evaluation.

\subsection{Evaluation Design}
\label{sec:v-final-design}

The orchestrator was invoked once, with a two-round repair budget
(\texttt{max\_rounds = 2}), over all 187 evaluation rows. For every
notebook, Round 1 is evaluated first. Round 2 is attempted only when Round 1
actually reached FixApplicator, came back \texttt{still\_failing}, and the
newly exposed error, once reclassified by the same classifier used for every
original row, is itself pip-repair-eligible. When Round 2 is attempted, it
runs the same RAGRepairAgent and FixApplicator components, targeting only
the newly exposed error, inside the container as Round 1 left it.

As explained in \Cref{sec:orchestrator}, Round 1 and Round 2 are not two
separate experiments: both conditions are read from the same run, so that
Ollama's non-zero sampling temperature cannot turn ordinary run-to-run
variation into an apparent Round-2 effect. The comparison that follows is
therefore between the outcome after Round 1 and the final outcome after up
to two rounds of the same execution, never between two independently
sampled runs.

Each of the 187 rows represents a notebook execution that had already been
observed to fail with a dependency error before this evaluation began.
Re-executing a notebook after a repair can expose a dependency error that
was not visible in the original dataset row, simply because execution now
reaches further into the notebook than it did before; \Cref{sec:v-final-interpretation}
returns to what this means for interpreting the results below.

\subsection{Primary Repair Results}
\label{sec:v-final-primary}

Full notebook recovery -- a notebook re-executing from beginning to end with
no error -- was not achieved for any evaluation notebook, in either round.
Round-1 repair success is 0/187 (0\%), and the final success rate after up
to two rounds is likewise 0/187 (0\%). Round 2 fixed no additional notebook
beyond what Round 1 had already fixed, for an absolute improvement of 0
percentage points. This is stated plainly and is not softened: no
evaluation notebook reached completely successful execution within the
maximum two-round repair budget defined for this system.

This does not mean that the repair actions themselves did nothing.
\Cref{sec:v-final-targeted} reports what actually happened at the level of
individual repair attempts, which is a different, and considerably more
informative, question than whether a notebook was fully recovered.

\subsection{Targeted Dependency-Error Resolution}
\label{sec:v-final-targeted}

Full notebook recovery requires every dependency problem in a notebook to be
absent after re-execution. A narrower, attempt-level question is whether the
one dependency error a given repair attempt specifically targeted was itself
removed -- regardless of whether the notebook went on to fail on something
else. A repair attempt counts as resolving its targeted error when the
notebook becomes \texttt{fixed}, or when it remains \texttt{still\_failing}
but the original error is gone and the conservative
\texttt{same\_as\_original\_error} check (\Cref{chap:methodology}) confirms
that the notebook progressed to a genuinely new, different error.

Across the 73 real FixApplicator attempts made in this evaluation (61 in
Round 1, 12 in Round 2), 62 resolved the error they targeted: a targeted
dependency-error resolution rate of 84.9\% (62/73). This is an attempt-level
figure, not a notebook-level one, and it should not be confused with the
0\% full-recovery figure above. Although no notebook completed successfully
within two rounds, the large majority of individual repair attempts did
remove the specific dependency error they addressed; re-execution then
frequently exposed a subsequent failure, which is exactly the pattern
already observed at pilot scale in \Cref{sec:v-key-observation} and traced
through one real record in \Cref{sec:m-worked-example}.

\subsection{Abstention: Round 1, Round 2, and Final State}
\label{sec:v-final-abstention}

A large share of the 187 evaluation rows never reached a repair attempt at
all, because RAGRepairAgent deliberately declined to propose one. This is
abstention, and it is a distinct outcome from an infrastructure fault or a
failed repair: the system does not generate an unsupported fix when it
lacks grounded package information, rather than guessing.

In Round 1, 126 of the 187 rows (67.4\%, the Round-1 abstention rate) were
abstained on. Every one of these 126 cases has the same cause: PyPI
retrieval returned \texttt{mapping\_unknown}, meaning RAGRepairAgent's
static import-to-distribution mapping (\Cref{chap:methodology}) could not
resolve the failing import to a supported PyPI package. In all 126 cases,
the language model was never called -- the abstention happens one layer
before it, at retrieval.

Round 2 does not simply repeat this figure. Of the 187 rows, 48 had a real
Round-1 attempt that came back \texttt{still\_failing} on a newly exposed,
repair-eligible error, and were therefore eligible for a second round.
Table~\ref{tab:v-final-round2} shows what happened to those 48. Only 12
received an actual second repair proposal and a real second FixApplicator
attempt; the remaining 36 abstained again, and, as in Round 1, every one of
these 36 abstentions is a \texttt{mapping\_unknown} case -- the newly exposed
error's import name had no entry in the package mapping either.

\begin{table}[h]
  \centering
  \begin{tabular}{|p{4.0cm}|r|p{6.5cm}|}
    \hline
    \textbf{Stage} & \textbf{Count} & \textbf{Meaning} \\
    \hline
    Round-2 eligible & 48 & Round 1 reached FixApplicator, came back \texttt{still\_failing}, and the newly exposed error is itself repair-eligible. \\
    \hline
    Round-2 proposal generated & 12 & RAGRepairAgent reached the language model and produced a proposal for the newly exposed error. \\
    \hline
    Round-2 abstained (\texttt{mapping\_unknown}) & 36 & Retrieval could not resolve the newly exposed error's import name; the model was never called. \\
    \hline
    Round-2 FixApplicator attempts & 12 & Every generated proposal was applied and the notebook re-executed. \\
    \hline
    Round-2 fixed & 0 & No second attempt produced a completely clean result. \\
    \hline
    Round-2 still failing & 12 & Every attempted second repair left the notebook failing on some error. \\
    \hline
  \end{tabular}
  \caption{What happened to the 48 Round-2-eligible notebooks. Round-2
  eligibility does not by itself mean a second repair was applied: for
  three of every four eligible notebooks, package resolution abstained
  before RAGRepairAgent could propose anything.}
  \label{tab:v-final-round2}
\end{table}

Combining both rounds, 162 of the 187 evaluation rows (86.6\%) ended in
abstention as their final classification: the 126 Round-1 abstentions plus
the 36 Round-2 abstentions above. This thesis reports these two figures
separately, as the Round-1 abstention rate (67.4\%) and the final-state
abstention rate (86.6\%), rather than under one undifferentiated "abstention
rate," because they answer different questions: the first is how often the
very first repair attempt could not proceed past package resolution, and the
second is how often a notebook ended the entire two-round process without
ever receiving a proposal. Every one of the 162 final-state abstentions,
in both rounds, has the same underlying cause -- \texttt{mapping\_unknown} --
which points at package-mapping coverage, not at the language model or at
the grounding logic, as the dominant reason a repair attempt never happened.

\subsection{Proposal Quality and Proposal Coverage}
\label{sec:v-final-proposal}

RAGRepairAgent was invoked, across both rounds, 235 times (187 in Round 1,
48 in Round 2). Of those invocations, 73 reached the point of actually
calling the language model; the remaining 162 abstained beforehand for the
\texttt{mapping\_unknown} reason described above. Of the 73 responses the
model produced, all 73 passed schema validation, and all 73 of those also
passed grounding validation against the retrieved PyPI evidence. In other
words, whenever RAGRepairAgent's own retrieval step supplied it with a
resolved candidate distribution, the model's proposal was schema-valid and
grounded without exception: a proposal validity rate of 100\% (73/73), a
grounding pass rate of 100\% (73/73 schema-valid proposals), and a grounded
proposal rate among LLM invocations of 100\% (73/73).

These three figures are conditional on the language model having actually
been called, and none of them should be read as a statement about overall
proposal coverage. Counting every one of the 235 repair-agent invocations,
including the 162 that abstained before reaching the model, only 73 ended
in a grounded proposal: an overall grounded proposal coverage rate of 31.1\%
(73/235). The contrast between these two figures is central to interpreting
this system: its output is reliable whenever it produces one at all, but the
static package-mapping coverage that decides whether the model is called in
the first place is the binding constraint on how often it gets that chance.

\subsection{Infrastructure Failures}
\label{sec:v-final-infra}

Two of the 187 evaluation rows (1.1\%) ended in an infrastructure failure --
an environment or tooling fault, unrelated to repair quality, that prevented
a fair trial of the method (\Cref{chap:methodology}). These are reported
separately from abstentions, method failures, and \texttt{still\_failing}
outcomes, and are not counted against the repair method itself, following
the same infrastructure-versus-method distinction already applied throughout
the pilots earlier in this chapter.

\subsection{Final Outcome Breakdown}
\label{sec:v-final-breakdown}

Table~\ref{tab:v-final-outcomes} summarises the final classification of
every one of the 187 evaluation rows, combining every figure reported above
into one place.

\begin{table}[h]
  \centering
  \begin{tabular}{|l|r|r|}
    \hline
    \textbf{Final category} & \textbf{Count} & \textbf{Share of 187} \\
    \hline
    \texttt{fixed} & 0 & 0.0\% \\
    \hline
    \texttt{still\_failing} & 22 & 11.8\% \\
    \hline
    \texttt{abstained} (126 Round 1 + 36 Round 2) & 162 & 86.6\% \\
    \hline
    \texttt{infrastructure\_failure} & 2 & 1.1\% \\
    \hline
    \texttt{method\_failure} & 1 & 0.5\% \\
    \hline
    \texttt{excluded} & 0 & 0.0\% \\
    \hline
    \textbf{Total} & \textbf{187} & \textbf{100\%} \\
    \hline
  \end{tabular}
  \caption{Final classification of all 187 evaluation rows, after up to two
  repair rounds.}
  \label{tab:v-final-outcomes}
\end{table}

\subsection{Results by Subtype}
\label{sec:v-final-subtype}

Split by the subtype recorded during classification (\Cref{chap:methodology}),
full notebook recovery is 0/169 (0\%) for \texttt{missing\_package} and 0/18
(0\%) for \texttt{wrong\_version}, the two subtypes within repair scope
(\Cref{chap:problem-dataset}). This full-recovery figure should be read
together with, not instead of, the attempt-level targeted-error resolution
rate reported in \Cref{sec:v-final-targeted}: a 0\% full-recovery rate for a
subtype does not mean individual repairs of that subtype's errors were
ineffective, only that no notebook containing one was left completely
error-free within the two-round budget.

\subsection{Manual Validation of the Classifier}
\label{sec:v-final-classifier}

To check ErrorClassifier's output against ground truth rather than only
against its own internal logic, a stratified sample of 49 records was frozen
before this evaluation was run and labelled by hand: all 10 records the
classifier had assigned \texttt{system\_library} and all 4 it had assigned
\texttt{mapping\_unknown} (the two subtypes excluded from repair scope), plus
20 of the records it had assigned \texttt{missing\_package} and 15 it had
assigned \texttt{wrong\_version}, drawn primarily from the evaluation split.

On the binary scope decision (usable versus excluded), the classifier scored
33 true positives, 2 false positives, 0 false negatives, and 14 true
negatives against the manual labels, for an accuracy of 95.9\%, a precision
of 94.3\%, a recall of 100\%, and an F1 score of 97.1\%. On the four-way
subtype decision, 45 of the 49 sample rows matched the manual label exactly,
a naive accuracy of 91.8\%. This figure describes the classifier's accuracy
on this deliberately stratified sample, which oversamples the two rarer,
excluded subtypes; it is reported as sample-level accuracy, not as a
population-wide estimate, since doing the latter without an independent
estimate of each subtype's true share of the full 214-row dataset would mean
using the classifier's own output to validate itself.

\begin{table}[h]
  \centering
  \begin{tabular}{|l|r|r|r|}
    \hline
    \textbf{Subtype} & \textbf{Precision} & \textbf{Recall} & \textbf{F1} \\
    \hline
    \texttt{missing\_package} & 0.800 & 1.000 & 0.889 \\
    \hline
    \texttt{wrong\_version} & 1.000 & 0.882 & 0.938 \\
    \hline
    \texttt{system\_library} & 1.000 & 1.000 & 1.000 \\
    \hline
    \texttt{mapping\_unknown} & 1.000 & 0.667 & 0.800 \\
    \hline
  \end{tabular}
  \caption{Per-subtype precision, recall, and F1 on the 49-record manual
  validation sample.}
  \label{tab:v-final-classifier}
\end{table}

Finally, the extracted failing-module name matched the manually confirmed
name exactly for all 49 sample rows (100\%). Extraction accuracy is reported
as an exact-match rate rather than as precision/recall/F1, since the space
of possible module names is unbounded and a confusion matrix over it would
not be meaningful.

\subsection{Manual Validation of PyPI Distribution Resolution}
\label{sec:v-final-pypi}

A separate, small diagnostic sample of 20 import names -- covering every
name already present in the package mapping and the most frequently
exercised names that are not -- was manually checked against the real PyPI
registry. Of the 20, 17 were names a human reviewer could confidently
resolve to one specific distribution or confirm as unresolvable; the
remaining 3 were left unresolved intentionally, because no single correct
distribution could be determined by hand either. Against those 17
manually resolvable cases, RAGRepairAgent's retrieval step resolved the
correct distribution in 7 (41.2\%).

This is a small, diagnostic sample, not a population-wide measurement of
package-resolution accuracy, and it is not presented as one. It is,
however, consistent with the abstention results in
\Cref{sec:v-final-abstention}: a package-mapping mechanism that resolves
fewer than half of a manually checked, deliberately diagnostic sample of
import names is the kind of coverage gap that would be expected to produce
a high \texttt{mapping\_unknown} abstention rate at full evaluation scale,
which is what was observed. This is offered as a consistent interpretation
of two independent pieces of evidence, not as proof that one caused the
other.

\subsection{Interpreting the Result: Why Zero Full Recoveries Does Not Mean Zero Repairs Worked}
\label{sec:v-final-interpretation}

Taken together, the results above show a system that behaves reliably at
the one stage that could be checked directly against ground truth --
every proposal the language model actually produced was schema-valid and
grounded -- but whose end-to-end repair coverage is limited by two separate
factors well before that stage is reached. Package-mapping coverage decides
whether RAGRepairAgent gets to call the model at all: on this evaluation
split, it did so for less than a third of all repair-agent invocations,
and abstained the rest of the time rather than guess a package name it
could not ground. Cascading dependency errors decide whether a notebook
that does receive a repair becomes fully runnable: fixing one import
frequently exposes a second, previously hidden one, and the bounded,
two-round repair budget defined for this system was not always enough to
clear every such error in the same notebook.

Neither factor is a defect in a single component. RAGRepairAgent's
abstention on \texttt{mapping\_unknown} is the deliberate, conservative
behaviour described in \Cref{chap:methodology}: refusing to act on
insufficient evidence rather than producing an unsupported suggestion.
FixApplicator's classification of a notebook as \texttt{still\_failing} on a
different error is likewise working exactly as designed
(\Cref{sec:v-key-observation}). The 0\% full-recovery result for this
evaluation split is the joint outcome of these two conservative, individually
correct behaviours operating within a bounded repair budget, not evidence
that the language model failed to produce useful repairs: it did, in 84.9\%
of the attempts it was actually given.

The results do not support a stronger claim than this. No notebook in the
evaluation split is shown to have needed more than two dependency errors in
total, nor is it shown that every notebook needed exactly two; some
Round-1 \texttt{still\_failing} outcomes were not even eligible for Round 2,
because the newly exposed error fell outside repair scope. The only claim
these results support directly is the one stated in
\Cref{sec:v-final-primary}: no evaluation notebook achieved complete
successful execution within the two-round repair budget evaluated here.

\section{Human Evaluation of Explanation Quality: Study Design}
\label{sec:v-human-eval}

\Cref{sec:v-final-eval} showed that LLMExplainer produced a schema-valid
explanation for 185 of the 187 evaluation rows. Schema validity is a
necessary check, but it is not the same claim as objective~O1: that the
explanation is written in plain language a non-expert reader can actually
follow (\texttt{NFR2}). Whether an explanation is understandable,
sufficiently detailed, and useful is a judgment only a human reader can
make, not something an automated schema check can establish. This section
reports the design of a small human evaluation built to close that gap. As
with the final evaluation, the design is reported separately from its
results: at the time of writing, the study has been fully designed and its
materials prepared, but it has not yet been administered to any
participant, and no response data exists.

\subsection{Instrument: An Adapted Explanation Satisfaction Scale}
\label{sec:v-human-eval-instrument}

Rather than devising a new questionnaire, this evaluation adapts a
published instrument: the Explanation Satisfaction Scale of Hoffman
et al.~\cite{hoffman2023measures}, a seven-item, 1--5 Likert scale for
collecting a user's own judgment of an explanation after they have read
it, as distinct from an expert's a priori judgment of an explanation's
designed-in properties (for which the same paper proposes a separate
Explanation Goodness Checklist, not used here). Six of the seven original
items are retained, each adapted only in its domain reference -- for
example, "From the explanation, I know how the [tool] works" becomes
"From this explanation, I understand why the notebook failed." The
seventh original item, which asks whether the explanation "tells me how
to use" the system, is excluded: LLMExplainer is deliberately barred from
suggesting what to do about a failure, since deriving a repair is
RAGRepairAgent's separate responsibility
(\Cref{sec:m-explanation-role}), so that item does not correspond to
anything this component is meant to provide. The retained item closest to
the original scale's accuracy question -- "This explanation seems
correct, based on the error shown" -- is reported throughout as
\emph{perceived correctness}, never as technical, objective, or factual
correctness: a participant's rating reflects only how correct the
explanation appears to the reader, not a verification against the
notebook's true root cause. Objective correctness is not measured by this
study at all.

The alternative System Causability Scale of Holzinger
et al.~\cite{holzinger2020scs} was considered and not adopted. Its ten
items presuppose an interactive explanation interface with an adjustable
level of detail, comparison across multiple explanations, and references
to domain-specific artefacts such as clinical guidelines -- properties of
an interactive, expert-facing medical decision-support tool that do not
apply to a single, static, plain-language explanation of one dependency
error.

All six retained items use the same 1--5 scale (1 = strongly disagree,
2 = disagree, 3 = neither agree nor disagree, 4 = agree, 5 = strongly
agree) after every example, so ratings stay comparable across the study.

\subsection{Participants}

Participants need only general programming experience; familiarity with
Python and Jupyter Notebooks is desirable but not required, and no
machine-learning or explainable-AI expertise is required, since the
explanations are written in plain language for exactly this audience
(\texttt{NFR2}). A short background section records self-reported general
programming experience, Python familiarity, Jupyter Notebook familiarity,
and prior experience with dependency/package errors -- kept deliberately
brief, and collecting no name, email, or other identifying information.
The target is 15--20 participants, a convenience sample sized for a
descriptive Master's-thesis study, not for a population-level claim.

\subsection{Example Pool: Diversity-Constrained Stratified Sampling}
\label{sec:v-human-eval-pool}

Twelve examples are drawn from the frozen 187-row evaluation split, never
from the 13-row development split (part of which is already quoted
verbatim elsewhere in this thesis, \Cref{sec:m-explanation-example}), and
never from five evaluation-split notebook\_execution\_ids already used in
earlier explanation-only development activity. Sampling is stratified by
subtype, but not forced into an equal split -- a design choice best
described as \emph{diversity-constrained stratified sampling}: the number
of examples drawn per subtype is bounded by how many genuinely distinct
error signatures exist in the eligible population, not fixed in advance.
Checking the eligible \texttt{wrong\_version} population directly showed
only three distinct (failing module, error message) signatures across all
18 eligible rows -- most of its records share one of two repeated numpy or
SciPy compatibility messages, which would produce near-duplicate
explanation text if sampled further. All three distinct signatures are
included, once each. The remaining nine pool positions are drawn from
\texttt{missing\_package}, whose eligible population supports far more
diversity (41 distinct failing modules across 162 eligible rows), giving
nine examples with nine different failing modules and no repeated
content. The result is \textbf{9 missing\_package and 3 wrong\_version}
examples, not a 6/6 split, and this composition is reported as such rather
than described as balanced. Because the pool contains nine distinct
\texttt{missing\_package} cases but only three distinct
\texttt{wrong\_version} cases, any subtype-level comparison in the
resulting data is explicitly descriptive: participants may contribute
several ratings of the three \texttt{wrong\_version} cases, but that does
not make them equivalent in diversity to nine distinct
\texttt{wrong\_version} explanations.

Selection used a fixed random seed (42) over alphabetically sorted
candidate lists, never consulted repair outcome, Round-2 status, or any
other automated-evaluation result, and is fully reproducible from the
frozen final-evaluation trace (\texttt{scripts/select\_human\_evaluation\_pool.py}).

Each participant rates 6 of the 12 pooled examples, assigned through four
fixed survey variants (\Cref{tab:v-human-eval-variants}) constructed so
that every one of the 12 examples appears in exactly two of the four
variants -- balanced exposure without needing a shared "anchor" example.
Participants are assigned variants by simple round robin as they are
recruited, and the presentation order of the 6 examples within a
participant's variant is randomized.

\begin{table}[h]
  \centering
  \begin{tabular}{|l|p{9.5cm}|}
    \hline
    \textbf{Variant} & \textbf{Notebook execution IDs (6 examples)} \\
    \hline
    A & 62, 79, 94, 118, 172, 163 \\
    \hline
    B & 62, 79, 94, 211, 252, 175 \\
    \hline
    C & 118, 211, 282, 398, 197, 163 \\
    \hline
    D & 172, 252, 282, 398, 175, 197 \\
    \hline
  \end{tabular}
  \caption{The four balanced survey variants. Every one of the 12 pooled
  examples appears in exactly two variants.}
  \label{tab:v-human-eval-variants}
\end{table}

\subsection{Procedure}

For each assigned example, a participant sees only the original error type
and error message, and the generated explanation's \texttt{summary},
\texttt{root\_cause}, \texttt{evidence}, and \texttt{limitations} fields,
rendered as plain text -- never as raw JSON, and never edited from what
the model actually produced. Nothing else is shown: not the repository
name, not any source code, not \texttt{explanation\_confidence}, and not
whether the notebook was ever repaired, whether a Round~2 attempt fired,
or any other outcome from the final evaluation. Outcome information is
withheld specifically because it could bias a participant's judgment of
the explanation itself, which is the only thing this study measures. The
same six items (\Cref{sec:v-human-eval-instrument}) are asked after every
example, in the same order, followed, once, by two optional open-text
questions asking what was unclear or missing and what would make the
explanations more helpful -- supplementary qualitative feedback, not part
of the validated scale.

\subsection{Planned Analysis}

For each of the six items, the plan is to report the number of responses,
median, interquartile range, the full 1--5 frequency distribution, and the
percentage rating 4 or 5, computed overall, by subtype, and per example.
No inferential significance test is planned, and no claim beyond this
specific sample of participants and examples is intended. Cronbach's
alpha is planned as a \emph{supplementary internal-consistency} statistic
only: it would describe whether the six adapted items tend to move
together in the collected data, not whether different participants agree
with each other (that is a separate question this design does not
address), and not whether the six items in fact measure one single
construct. A high alpha would not by itself justify collapsing the six
items into one composite "explanation quality score," and no such
composite is computed by default; item-level results remain the primary
output regardless of the alpha value obtained.

\subsection{Limitations of This Design}

Three limitations are stated plainly rather than discovered later. First,
participants judge a rendered error message and explanation, not the
original notebook and its surrounding code, so a genuinely different
impression might emerge with fuller context. Second, the diversity
constraint on \texttt{wrong\_version} means the study can describe, but
not statistically compare, perceived quality across the two subtypes with
equal confidence. Third, a convenience sample of 15--20 participants
supports a descriptive characterization of how this specific group
perceived this specific system's explanations, not a claim generalizable
to programmers at large.

\section{Chapter Summary}
\label{sec:v-summary}

This chapter reported several different kinds of evidence, in increasing order
of scale. The automated test suite (\Cref{sec:v-testing}) shows that every
implemented component -- error classification, explanation, PyPI-grounded
repair proposal generation, fix application with re-execution, and result
logging into both the SQLite benchmark table and the FAIR Jupyter knowledge
graph -- behaves as designed against controlled, mocked inputs. The
component-by-component pilots (\Cref{sec:v-i4,sec:v-i5,sec:v-i6}) and the
small orchestrator pilot (\Cref{sec:v-i7}) then showed that each of those
components, and the orchestrator that connects them, also behaves correctly
against its relevant real external system -- a real Ollama model, the live
PyPI registry, a real Docker daemon, real GitHub repositories, and the real
Morph-KGC mapping engine -- and surfaced concrete, fixable problems (a
timeout set too low, two missing package mappings, three
environment-specific bugs) that only a real pilot could have found.

\Cref{sec:v-final-eval} then reported the thesis's final evaluation: the
same, unmodified repair layer run once over the complete 187-row evaluation
split. No evaluation notebook reached full recovery within the two-round
repair budget (0/187), but 84.9\% of the 73 real repair attempts made
resolved the dependency error they specifically targeted, and every one of
the 73 proposals the language model actually produced was schema-valid and
grounded. Repair coverage was limited well before that point: a
\texttt{mapping\_unknown} package-resolution gap accounted for all 162
final-state abstentions (126 in Round 1, a further 36 in Round 2), leaving
only 73 of 235 repair-agent invocations (31.1\%) reaching the language model
at all. Two rows (1.1\%) ended in an infrastructure failure, kept separate
from these method-level figures throughout.

What this evaluation does not show is a repair success rate for the whole
214-row dataset, since 14 rows fall outside pip-only repair scope by design
(\Cref{chap:problem-dataset}), or a knowledge graph enriched at the same
scale as the evaluation run: the \texttt{repair\_attempts} table now holds
one row per Round-1 and Round-2 attempt from the full evaluation split, but
running the existing RML mapping over that larger table and a baseline
comparison against a non-grounded repair agent remain future work, as
\Cref{chap:future-work} describes. \Cref{sec:v-human-eval} reports one
further piece of evidence this chapter adds: a human evaluation of
explanation quality, whose design is now finalized and whose materials are
prepared, but which has not yet been administered to any participant.
